% Chapter 5
\chapter{Prediction Models and Experimental Evaluation}
\label{Chapter5}
\lhead{Chapter 5. \emph{Prediction Models and Experimental Evaluation}}

This chapter describes the recurrent predictors used in Magnavis, the offline benchmark that compares six detection schemes, and the principal findings on short- and long-duration ground truth.

\section{Forecasting Models}
\subsection{Shared preprocessing}
Each one-second sample of resultant magnetic magnitude is scaled and paired with cyclic time-of-day features (sine and cosine of the hour angle). A sliding window of length $W$ (default 15 in the benchmark suite) forms the input; the target is the next magnitude. During inference the model emits a one-step-ahead forecast that is compared with the measured value.

\subsection{Pretrained GRU and LSTM}
Both families use a recurrent core followed by dense layers. Checkpoints are stored under \texttt{models/} as \texttt{gru\_pretrained\_OBS2\_*.keras} and \texttt{lstm\_pretrained\_OBS2\_*.keras} with matching scaler files. Training on a December 2025 OBS2 CSV archive is performed offline via \texttt{train\_gru\_pretrained.py} and \texttt{train\_lstm\_pretrained.py}. At runtime the headless application loads these weights and runs in \emph{predict-only} mode when no on-session training window is configured.

\subsection{Attention bidirectional LSTM (fresh training)}
For comparison, an attention bidirectional LSTM is trained inside the benchmark on the first 62 minutes of each evaluation CSV, then used to forecast the remainder. This represents an on-site adaptation strategy rather than a long-archive pretrain.

\subsection{Statistical baselines}
Three non-neural baselines forecast the field by smoothing the observed series (exponential moving average, median filter, and Savitzky--Golay). Anomalies are still detected on the residual between observation and smoother output, so all schemes share the same downstream detector.

\section{Benchmark Protocol}
The evaluation bundle \texttt{k\_recall\_curves\_zero\_hist\_20260522\_170707} implements a \textbf{zero-historic, predict-only} protocol:
\begin{itemize}
    \item No leading training minutes inside the app session; pretrained weights are used as shipped.
    \item Historic context set to zero; skip-initial minutes set to zero unless noted otherwise.
    \item Detectors operate on OBS2 sensors 1--3 with aligned timestamps.
    \item Residual $e_t = |B_t - \hat B_t|$; exponential moving average yields $\mu_t$ and $\sigma_t$; flag if $e_t > \mu_t + k\sigma_t$ for $k \in \{1,\ldots,5\}$.
    \item Point-level metrics (TP, FP, TN, FN, recall, precision, F1) are computed on the same second grid for every model.
\end{itemize}

\subsection{Datasets and ground-truth layouts}
\begin{table}[htbp]
\centering
\caption{Benchmark scenarios and class balance.}
\label{tab:benchmark_scenarios}
\begin{tabular}{|l|l|r|r|}
\hline
\textbf{Key} & \textbf{Description} & \textbf{Evaluated points} & \textbf{GT-positive rate} \\
\hline
Feb-13 & Short shot-duration windows & 4,711 & $\approx 4.4\%$ \\
Apr-27 & Long sustained-offset windows & 8,825 & $\approx 94.5\%$ \\
Synthetic & Flat magnetics with long-window labels & 10,800 & $\approx 50.0\%$ \\
\hline
\end{tabular}
\end{table}

Additional runs compare open-loop and closed-loop autoregressive inference on Apr-27 for GRU and LSTM (\texttt{PREDICTOR\_AR\_CLOSED\_LOOP=1}).

\section{Results}
\subsection{Short-duration ground truth (Feb-13)}
On sparse anomalies, statistical smoothers achieve high recall at low $k$ but lose recall as $k$ increases. Pretrained GRU and LSTM remain competitive through $k=4$ (recall about 85--88\%). Attention Bi-LSTM trained on the leading window reaches 92\% recall at $k=4$ with higher precision, showing that local adaptation helps when anomalies are brief and rare.

\subsection{Long-duration ground truth (Apr-27)}
This scenario exposes the largest differences between models. Simple smoothers collapse above $k=2$ because they cannot track sustained offsets. Pretrained LSTM holds roughly constant 85\% recall for $k=2$--$5$. Pretrained GRU matches LSTM at $k=3$ (85\% recall) but drops to 63\% for $k \ge 4$, losing about 1,800 true anomaly seconds between $k=3$ and $k=4$. Precision remains high ($\approx 94\%$); the failure is missed interior seconds, not excessive false alarms.

Figure~\ref{fig:k_recall_five} summarises recall versus $k$ for all five scenario panels. Figure~\ref{fig:apr27_open_closed} shows that closed-loop GRU inference does not restore long-GT recall at $k=4$.

\begin{figure}[htbp]
\centering
\includegraphics[width=\textwidth]{k_recall_curves_five_scenarios.png}
\caption{Recall versus detector multiplier $k$ for six schemes across five scenario panels (zero-historic benchmark).}
\label{fig:k_recall_five}
\end{figure}

\begin{figure}[htbp]
\centering
\includegraphics[width=0.85\textwidth]{k_recall_apr27_open_vs_closed_loop.png}
\caption{Apr-27 long ground truth: open-loop versus closed-loop GRU and LSTM at $k=4$.}
\label{fig:apr27_open_closed}
\end{figure}

\subsection{Mechanism for the GRU cliff}
When a predictor tracks a sustained level shift, the residual $|B-\hat B|$ shrinks in the interior of a long labelled interval. Raising $k$ raises the threshold; interior seconds then fall below the threshold even though the ground truth remains positive. LSTM checkpoints used here maintain larger interior residuals on Apr-27; GRU checkpoints track the plateau more tightly. A retraining ablation with absolute targets restores flat recall across $k$, confirming that the effect is linked to predictor behaviour rather than to the labelling file alone.

\subsection{Synthetic long-window layout}
With flat synthetic magnetics and long-window labels, GRU recall stays near 100\% for all $k$. The Apr-27 cliff therefore requires both the long label layout and realistic field shape during the event, not the label geometry alone.

\subsection{Deployment guidance}
For long campaigns with sustained anomalies, pretrained LSTM at $k \approx 3$ is the most stable choice in this study. GRU at $k \ge 4$ on long ground truth is not recommended without retraining or threshold tuning. Short campaigns may use $k=4$ with either recurrent model. Full point-level counts appear in Tables~\ref{tab:anomaly_k_recall_zero_hist_I} and~\ref{tab:anomaly_k_recall_zero_hist_II}.

% Requires: \usepackage{multirow}
% Model column: \multirow spans all scenario×k rows; scenario rules use
% \cline{2-12} so horizontal lines do not cut through the merged model cell.
\begin{table*}[t]
\centering
\footnotesize
\caption{Point-level anomaly detection across six schemes, three open-loop scenarios (short / long / synthetic GT), and GRU/LSTM closed-loop on Apr-27 long GT (\texttt{PREDICTOR\_AR\_CLOSED\_LOOP=1}), for $k \in \{1,\ldots,5\}$ (zero-historic, predict-only benchmark). Short Duration: 4{,}711 s; Long Duration: 8{,}825 s; Synthetic: 10{,}800 s. Part~I: EWMA, Median, Savitzky--Golay, and Attention Bi-LSTM.}
\label{tab:anomaly_k_recall_zero_hist_I}
\resizebox{\textwidth}{!}{
\begin{tabular}{|l|l|c|r|r|r|r|c|c|c|c|c|}
\hline
\textbf{Model} & \textbf{Scenario} & $\mathbf{K}$ & \textbf{TP} & \textbf{FP} & \textbf{TN} & \textbf{FN} & \textbf{Recall} & \textbf{Precision} & \textbf{F1} & \textbf{Specificity} & \textbf{Accuracy} \\
\hline
\multirow{15}{*}{EWMA} & \multirow{5}{*}{Short Duration Anomaly} & 1 & 207 & 4450 & 54 & 0 & 1.0000 & 0.0444 & 0.0851 & 0.0120 & 0.0554 \\
 &  & 2 & 200 & 1201 & 3303 & 7 & 0.9662 & 0.1428 & 0.2488 & 0.7333 & 0.7436 \\
 &  & 3 & 195 & 149 & 4355 & 12 & 0.9420 & 0.5669 & 0.7078 & 0.9669 & 0.9658 \\
 &  & 4 & 188 & 40 & 4464 & 19 & 0.9082 & 0.8246 & 0.8644 & 0.9911 & 0.9875 \\
 &  & 5 & 180 & 29 & 4475 & 27 & 0.8696 & 0.8612 & 0.8654 & 0.9936 & 0.9881 \\
\cline{2-12}
 & \multirow{5}{*}{Long Duration Anomaly} & 1 & 8293 & 481 & 0 & 51 & 0.9939 & 0.9452 & 0.9689 & 0.0000 & 0.9397 \\
 &  & 2 & 2174 & 90 & 391 & 6170 & 0.2605 & 0.9602 & 0.4099 & 0.8129 & 0.2907 \\
 &  & 3 & 225 & 44 & 437 & 8119 & 0.0270 & 0.8364 & 0.0522 & 0.9085 & 0.0750 \\
 &  & 4 & 39 & 36 & 445 & 8305 & 0.0047 & 0.5200 & 0.0093 & 0.9252 & 0.0548 \\
 &  & 5 & 19 & 34 & 447 & 8325 & 0.0023 & 0.3585 & 0.0045 & 0.9293 & 0.0528 \\
\cline{2-12}
 & \multirow{5}{*}{Synthetic Data} & 1 & 87 & 5366 & 31 & 5316 & 0.0161 & 0.0160 & 0.0160 & 0.0057 & 0.0109 \\
 &  & 2 & 75 & 4958 & 439 & 5328 & 0.0139 & 0.0149 & 0.0144 & 0.0813 & 0.0476 \\
 &  & 3 & 68 & 4525 & 872 & 5335 & 0.0126 & 0.0148 & 0.0136 & 0.1616 & 0.0870 \\
 &  & 4 & 63 & 4495 & 902 & 5340 & 0.0117 & 0.0138 & 0.0126 & 0.1671 & 0.0894 \\
 &  & 5 & 47 & 4302 & 1095 & 5356 & 0.0087 & 0.0108 & 0.0096 & 0.2029 & 0.1057 \\
\hline
\multirow{15}{*}{Median} & \multirow{5}{*}{Short Duration Anomaly} & 1 & 207 & 4453 & 51 & 0 & 1.0000 & 0.0444 & 0.0851 & 0.0113 & 0.0548 \\
 &  & 2 & 205 & 1576 & 2928 & 2 & 0.9903 & 0.1151 & 0.2062 & 0.6501 & 0.6650 \\
 &  & 3 & 204 & 231 & 4273 & 3 & 0.9855 & 0.4690 & 0.6355 & 0.9487 & 0.9503 \\
 &  & 4 & 201 & 71 & 4433 & 6 & 0.9710 & 0.7390 & 0.8392 & 0.9842 & 0.9837 \\
 &  & 5 & 196 & 46 & 4458 & 11 & 0.9469 & 0.8099 & 0.8731 & 0.9898 & 0.9879 \\
\cline{2-12}
 & \multirow{5}{*}{Long Duration Anomaly} & 1 & 8295 & 479 & 2 & 49 & 0.9941 & 0.9454 & 0.9692 & 0.0042 & 0.9402 \\
 &  & 2 & 2603 & 121 & 360 & 5741 & 0.3120 & 0.9556 & 0.4704 & 0.7484 & 0.3358 \\
 &  & 3 & 272 & 58 & 423 & 8072 & 0.0326 & 0.8242 & 0.0627 & 0.8794 & 0.0788 \\
 &  & 4 & 52 & 50 & 431 & 8292 & 0.0062 & 0.5098 & 0.0123 & 0.8960 & 0.0547 \\
 &  & 5 & 34 & 49 & 432 & 8310 & 0.0041 & 0.4096 & 0.0081 & 0.8981 & 0.0528 \\
\cline{2-12}
 & \multirow{5}{*}{Synthetic Data} & 1 & 51 & 5358 & 39 & 5352 & 0.0094 & 0.0094 & 0.0094 & 0.0072 & 0.0083 \\
 &  & 2 & 51 & 5045 & 352 & 5352 & 0.0094 & 0.0100 & 0.0097 & 0.0652 & 0.0373 \\
 &  & 3 & 51 & 4562 & 835 & 5352 & 0.0094 & 0.0111 & 0.0102 & 0.1547 & 0.0820 \\
 &  & 4 & 51 & 4503 & 894 & 5352 & 0.0094 & 0.0112 & 0.0102 & 0.1656 & 0.0875 \\
 &  & 5 & 51 & 4307 & 1090 & 5352 & 0.0094 & 0.0117 & 0.0104 & 0.2020 & 0.1056 \\
\hline
\multirow{15}{*}{Savitzky--Golay} & \multirow{5}{*}{Short Duration Anomaly} & 1 & 207 & 4324 & 180 & 0 & 1.0000 & 0.0457 & 0.0874 & 0.0400 & 0.0821 \\
 &  & 2 & 206 & 1142 & 3362 & 1 & 0.9952 & 0.1528 & 0.2650 & 0.7464 & 0.7574 \\
 &  & 3 & 205 & 184 & 4320 & 2 & 0.9903 & 0.5270 & 0.6879 & 0.9591 & 0.9605 \\
 &  & 4 & 203 & 77 & 4427 & 4 & 0.9807 & 0.7250 & 0.8337 & 0.9829 & 0.9828 \\
 &  & 5 & 195 & 52 & 4452 & 12 & 0.9420 & 0.7895 & 0.8590 & 0.9885 & 0.9864 \\
\cline{2-12}
 & \multirow{5}{*}{Long Duration Anomaly} & 1 & 8083 & 474 & 7 & 261 & 0.9687 & 0.9446 & 0.9565 & 0.0146 & 0.9167 \\
 &  & 2 & 2055 & 125 & 356 & 6289 & 0.2463 & 0.9427 & 0.3905 & 0.7401 & 0.2732 \\
 &  & 3 & 265 & 73 & 408 & 8079 & 0.0318 & 0.7840 & 0.0610 & 0.8482 & 0.0763 \\
 &  & 4 & 79 & 70 & 411 & 8265 & 0.0095 & 0.5302 & 0.0186 & 0.8545 & 0.0555 \\
 &  & 5 & 57 & 70 & 411 & 8287 & 0.0068 & 0.4488 & 0.0135 & 0.8545 & 0.0530 \\
\cline{2-12}
 & \multirow{5}{*}{Synthetic Data} & 1 & 93 & 5312 & 85 & 5310 & 0.0172 & 0.0172 & 0.0172 & 0.0157 & 0.0165 \\
 &  & 2 & 93 & 4834 & 563 & 5310 & 0.0172 & 0.0189 & 0.0180 & 0.1043 & 0.0607 \\
 &  & 3 & 93 & 4553 & 844 & 5310 & 0.0172 & 0.0200 & 0.0185 & 0.1564 & 0.0868 \\
 &  & 4 & 93 & 4508 & 889 & 5310 & 0.0172 & 0.0202 & 0.0186 & 0.1647 & 0.0909 \\
 &  & 5 & 93 & 4386 & 1011 & 5310 & 0.0172 & 0.0208 & 0.0188 & 0.1873 & 0.1022 \\
\hline
\multirow{15}{*}{Attention Bi-LSTM} & \multirow{5}{*}{Short Duration Anomaly} & 1 & 207 & 4322 & 182 & 0 & 1.0000 & 0.0457 & 0.0874 & 0.0404 & 0.0826 \\
 &  & 2 & 205 & 929 & 3575 & 2 & 0.9903 & 0.1808 & 0.3057 & 0.7937 & 0.8024 \\
 &  & 3 & 199 & 50 & 4454 & 8 & 0.9614 & 0.7992 & 0.8728 & 0.9889 & 0.9877 \\
 &  & 4 & 190 & 5 & 4499 & 17 & 0.9179 & 0.9744 & 0.9453 & 0.9989 & 0.9953 \\
 &  & 5 & 188 & 5 & 4499 & 19 & 0.9082 & 0.9741 & 0.9400 & 0.9989 & 0.9949 \\
\cline{2-12}
 & \multirow{5}{*}{Long Duration Anomaly} & 1 & 7419 & 476 & 5 & 925 & 0.8891 & 0.9397 & 0.9137 & 0.0104 & 0.8412 \\
 &  & 2 & 7096 & 425 & 56 & 1248 & 0.8504 & 0.9435 & 0.8945 & 0.1164 & 0.8104 \\
 &  & 3 & 5344 & 425 & 56 & 3000 & 0.6405 & 0.9263 & 0.7573 & 0.1164 & 0.6119 \\
 &  & 4 & 7066 & 421 & 60 & 1278 & 0.8468 & 0.9438 & 0.8927 & 0.1247 & 0.8075 \\
 &  & 5 & 7077 & 268 & 213 & 1267 & 0.8482 & 0.9635 & 0.9022 & 0.4428 & 0.8261 \\
\cline{2-12}
 & \multirow{5}{*}{Synthetic Data} & 1 & 5398 & 4676 & 721 & 5 & 0.9991 & 0.5358 & 0.6976 & 0.1336 & 0.5666 \\
 &  & 2 & 5398 & 4434 & 963 & 5 & 0.9991 & 0.5490 & 0.7086 & 0.1784 & 0.5890 \\
 &  & 3 & 5403 & 3157 & 2240 & 0 & 1.0000 & 0.6312 & 0.7739 & 0.4150 & 0.7077 \\
 &  & 4 & 3598 & 4497 & 900 & 1805 & 0.6659 & 0.4445 & 0.5331 & 0.1668 & 0.4165 \\
 &  & 5 & 3595 & 4495 & 902 & 1808 & 0.6654 & 0.4444 & 0.5329 & 0.1671 & 0.4164 \\
\hline
\end{tabular}
}
\end{table*}

\begin{table*}[p]
\centering
\footnotesize
\caption{Point-level anomaly detection across six schemes, three open-loop scenarios (short / long / synthetic GT), and GRU/LSTM closed-loop on Apr-27 long GT (\texttt{PREDICTOR\_AR\_CLOSED\_LOOP=1}), for $k \in \{1,\ldots,5\}$ (zero-historic, predict-only benchmark). Short Duration: 4{,}711 s; Long Duration: 8{,}825 s; Synthetic: 10{,}800 s. Part~II: GRU and LSTM (including Apr-27 closed-loop).}
\label{tab:anomaly_k_recall_zero_hist_II}
\resizebox{\textwidth}{!}{
\begin{tabular}{|l|l|c|r|r|r|r|c|c|c|c|c|}
\hline
\textbf{Model} & \textbf{Scenario} & $\mathbf{K}$ & \textbf{TP} & \textbf{FP} & \textbf{TN} & \textbf{FN} & \textbf{Recall} & \textbf{Precision} & \textbf{F1} & \textbf{Specificity} & \textbf{Accuracy} \\
\hline
\multirow{20}{*}{GRU} & \multirow{5}{*}{Short Duration Anomaly} & 1 & 202 & 3717 & 787 & 5 & 0.9758 & 0.0515 & 0.0979 & 0.1747 & 0.2099 \\
 &  & 2 & 191 & 1831 & 2673 & 16 & 0.9227 & 0.0945 & 0.1714 & 0.5935 & 0.6079 \\
 &  & 3 & 187 & 390 & 4114 & 20 & 0.9034 & 0.3241 & 0.4770 & 0.9134 & 0.9130 \\
 &  & 4 & 183 & 41 & 4463 & 24 & 0.8841 & 0.8170 & 0.8492 & 0.9909 & 0.9862 \\
 &  & 5 & 183 & 22 & 4482 & 24 & 0.8841 & 0.8927 & 0.8883 & 0.9951 & 0.9902 \\
\cline{2-12}
 & \multirow{5}{*}{Long Duration Anomaly} & 1 & 8205 & 479 & 2 & 139 & 0.9833 & 0.9448 & 0.9637 & 0.0042 & 0.9300 \\
 &  & 2 & 7291 & 432 & 49 & 1053 & 0.8738 & 0.9441 & 0.9076 & 0.1019 & 0.8317 \\
 &  & 3 & 7125 & 428 & 53 & 1219 & 0.8539 & 0.9433 & 0.8964 & 0.1102 & 0.8134 \\
 &  & 4 & 5281 & 349 & 132 & 3063 & 0.6329 & 0.9380 & 0.7558 & 0.2744 & 0.6134 \\
 &  & 5 & 5275 & 179 & 302 & 3069 & 0.6322 & 0.9672 & 0.7646 & 0.6279 & 0.6320 \\
\cline{2-12}
 & \multirow{5}{*}{Synthetic Data} & 1 & 5403 & 4571 & 826 & 0 & 1.0000 & 0.5417 & 0.7027 & 0.1530 & 0.5768 \\
 &  & 2 & 5403 & 4000 & 1397 & 0 & 1.0000 & 0.5746 & 0.7298 & 0.2588 & 0.6296 \\
 &  & 3 & 5403 & 3668 & 1729 & 0 & 1.0000 & 0.5956 & 0.7466 & 0.3204 & 0.6604 \\
 &  & 4 & 5403 & 2352 & 3045 & 0 & 1.0000 & 0.6967 & 0.8212 & 0.5642 & 0.7822 \\
 &  & 5 & 5403 & 1983 & 3414 & 0 & 1.0000 & 0.7315 & 0.8449 & 0.6326 & 0.8164 \\
\cline{2-12}
 & \multirow{5}{*}{GRU Closed-loop} & 1 & 8309 & 481 & 0 & 35 & 0.9958 & 0.9453 & 0.9699 & 0.0000 & 0.9415 \\
 &  & 2 & 8300 & 481 & 0 & 44 & 0.9947 & 0.9452 & 0.9693 & 0.0000 & 0.9405 \\
 &  & 3 & 7090 & 428 & 53 & 1254 & 0.8497 & 0.9431 & 0.8940 & 0.1102 & 0.8094 \\
 &  & 4 & 5386 & 388 & 93 & 2958 & 0.6455 & 0.9328 & 0.7630 & 0.1933 & 0.6208 \\
 &  & 5 & 4766 & 108 & 373 & 3578 & 0.5712 & 0.9778 & 0.7211 & 0.7755 & 0.5823 \\
\hline
\multirow{20}{*}{LSTM} & \multirow{5}{*}{Short Duration Anomaly} & 1 & 204 & 3804 & 700 & 3 & 0.9855 & 0.0509 & 0.0968 & 0.1554 & 0.1919 \\
 &  & 2 & 186 & 936 & 3568 & 21 & 0.8986 & 0.1658 & 0.2799 & 0.7922 & 0.7969 \\
 &  & 3 & 179 & 58 & 4446 & 28 & 0.8647 & 0.7553 & 0.8063 & 0.9871 & 0.9817 \\
 &  & 4 & 176 & 15 & 4489 & 31 & 0.8502 & 0.9215 & 0.8844 & 0.9967 & 0.9902 \\
 &  & 5 & 152 & 12 & 4492 & 55 & 0.7343 & 0.9268 & 0.8194 & 0.9973 & 0.9858 \\
\cline{2-12}
 & \multirow{5}{*}{Long Duration Anomaly} & 1 & 7079 & 430 & 51 & 1265 & 0.8484 & 0.9427 & 0.8931 & 0.1060 & 0.8079 \\
 &  & 2 & 7078 & 429 & 52 & 1266 & 0.8483 & 0.9429 & 0.8931 & 0.1081 & 0.8079 \\
 &  & 3 & 7078 & 429 & 52 & 1266 & 0.8483 & 0.9429 & 0.8931 & 0.1081 & 0.8079 \\
 &  & 4 & 7078 & 428 & 53 & 1266 & 0.8483 & 0.9430 & 0.8931 & 0.1102 & 0.8080 \\
 &  & 5 & 7078 & 428 & 53 & 1266 & 0.8483 & 0.9430 & 0.8931 & 0.1102 & 0.8080 \\
\cline{2-12}
 & \multirow{5}{*}{Synthetic Data} & 1 & 5403 & 4591 & 806 & 0 & 1.0000 & 0.5406 & 0.7018 & 0.1493 & 0.5749 \\
 &  & 2 & 5403 & 3877 & 1520 & 0 & 1.0000 & 0.5822 & 0.7360 & 0.2816 & 0.6410 \\
 &  & 3 & 5403 & 2392 & 3005 & 0 & 1.0000 & 0.6931 & 0.8188 & 0.5568 & 0.7785 \\
 &  & 4 & 5403 & 701 & 4696 & 0 & 1.0000 & 0.8852 & 0.9391 & 0.8701 & 0.9351 \\
 &  & 5 & 5403 & 88 & 5309 & 0 & 1.0000 & 0.9840 & 0.9919 & 0.9837 & 0.9919 \\
\cline{2-12}
 & \multirow{5}{*}{LSTM Closed-loop} & 1 & 8309 & 481 & 0 & 35 & 0.9958 & 0.9453 & 0.9699 & 0.0000 & 0.9415 \\
 &  & 2 & 7852 & 429 & 52 & 492 & 0.9410 & 0.9482 & 0.9446 & 0.1081 & 0.8956 \\
 &  & 3 & 7090 & 428 & 53 & 1254 & 0.8497 & 0.9431 & 0.8940 & 0.1102 & 0.8094 \\
 &  & 4 & 7079 & 428 & 53 & 1265 & 0.8484 & 0.9430 & 0.8932 & 0.1102 & 0.8082 \\
 &  & 5 & 4925 & 92 & 389 & 3419 & 0.5902 & 0.9817 & 0.7372 & 0.8087 & 0.6022 \\
\hline
\end{tabular}
}
\end{table*}
% 100 data cells (60 + 40)


\section{Direction of the Magnetic Perturbation}
When CSV or database rows provide three orthogonal components, the application subtracts a quiet-period baseline (median per component) and forms the perturbation vector $\Delta\mathbf{B}$. Azimuth and inclination follow standard spherical conventions; a unit vector is drawn in the 3D panel. If two observatories flag anomalies within a short interval, bearing lines can be intersected for a rough source location.

This step estimates the direction of $\Delta\mathbf{B}$, not necessarily the bearing to a compact source. Near-field geometry, baseline drift, and timing mismatch between components can bias the result. Chapter~\ref{Chapter7} lists improvements such as calibration runs and dipole inversion.

\section*{Summary}
Recurrent predictors outperform direct smoothing on long-duration ground truth. LSTM is more stable across threshold settings; GRU shows a sharp recall drop at $k \ge 4$ on Apr-27 because of tight plateau tracking. The benchmark scripts and result tables in the repository allow independent reproduction of these figures.
